\chapter{Restricciones y Metodología}
\label{cap:restricciones_metodologia}

El desarrollo de Luxe Core AI como TFG impone un marco de trabajo delimitado por restricciones técnicas, temporales y económicas que han moldeado las decisiones de diseño. Este capítulo las detalla y explica los \textit{trade-offs} asumidos.


\section{Restricciones del Proyecto}

\subsection{Restricción Temporal}

El límite de un curso académico es la restricción más inflexible del proyecto. Define de forma estricta el alcance del producto mínimo viable (MVP). La metodología adoptada prioriza la iteración rápida y las funcionalidades troncales operativas frente a la sobreingeniería: es mejor tener un sistema que funcione con pocos dispositivos que uno diseñado para mil que nunca llegue a ejecutarse.

\subsection{Restricción Económica}

El proyecto ha operado con un presupuesto de hardware deliberadamente reducido. El autor es estudiante sin ingresos, lo que ha llevado a seleccionar los dispositivos más económicos disponibles (como el ventilador FanLamp F8808, el modelo más barato de AliExpress con conectividad) y a reutilizar hardware existente (el ESP32 DevKit V1, originalmente adquirido para el puente HVAC, se empleó también como radio BLE programable para la ingeniería inversa del ventilador).

Se ha descartado cualquier compra de hardware auxiliar no estrictamente imprescindible: nada de dongles analizadores BLE, mandos universales o \textit{sniffers} de protocolos. Esta restricción no es una limitación del proyecto, sino una característica deliberada. Obliga a desarrollar competencias de ingeniería inversa que son transferibles a cualquier dispositivo IoT futuro, y muestra que el control del entorno doméstico con IA no tiene por qué requerir grandes inversiones.


\section{Decisiones Arquitectónicas}

\subsection{Evolución de la Infraestructura}

El desarrollo de modelos de lenguaje locales requiere capacidad de cómputo que un portátil de estudiante no siempre puede proporcionar. Durante la primera fase del proyecto, se utilizó la infraestructura de la Universidad de Córdoba (servidores con GPU RTX, accesibles por VPN y SSH) para prototipar la arquitectura de agentes. Esta fase fue útil para validar el enfoque, pero dependía de una conexión externa y de recursos que el autor no controlaba.

En la segunda fase, el sistema migró a un servidor doméstico dedicado (torre Ryzen APU, 16~GB RAM). Esta migración no fue un retroceso, sino la materialización del objetivo central del proyecto: eliminar cualquier dependencia externa. Hoy, el sistema funciona al 100\% en local, con modelos de lenguaje ejecutándose en Ollama sobre el propio servidor, sin necesidad de GPU dedicada ni de conexión a Internet.

Durante esta fase también se evaluaron plataformas \textit{open-source} de IoT como ESPHome, Tasmota y OpenMQTTGateway como posible base para las integraciones, pero se optó por un enfoque más ligero basado en \textit{drivers} Python específicos para cada dispositivo, priorizando la simplicidad y el control sobre el hardware.

\subsection{Paradigma Políglota}

La naturaleza dual de la domótica con IA (tiempo real para el control hardware, flexibilidad cognitiva para la toma de decisiones) hace inviable un sistema monolítico en un solo lenguaje. Se ha adoptado una arquitectura de tres capas con lenguajes especializados:

\begin{itemize}
    \item \textbf{Python 3.12} para los \textit{drivers} de dispositivos (Tuya, BLE, FanLamp) y el enrutador de modelos. Python ofrece un ecosistema maduro de librerías (\texttt{tinytuya}~\cite{tinytuya}, \texttt{bleak}~\cite{bleak}) y facilidad de prototipado.
    \item \textbf{Node.js 24} para el orquestador OpenClaw~\cite{openclaw}, que proporciona los canales de usuario (Telegram, \textit{dashboard} web) y la interfaz de agentes de IA.
    \item \textbf{C++20} para el firmware del ESP32, donde se necesita control determinista de memoria y baja latencia en la comunicación con el bus LIN. \textbf{[Pendiente de implementar.]}
\end{itemize}

Esta separación permite elegir el lenguaje óptimo para cada capa sin acoplar los tiempos de ejecución ni forzar compromisos de rendimiento. Herramientas de asistencia de código como OpenCode o Claude Code pueden ejecutarse directamente en el servidor a través de la sesión SSH, beneficiándose de su capacidad de cómputo sin la latencia que implicaría un IDE remoto.

\subsection{Evolución de la Selección de Modelos}
\label{sec:migracion_qwen}

La configuración actual de modelos es el resultado de un proceso iterativo de evaluación empírica sobre el hardware real del proyecto (Ryzen 5, 16~GB RAM, sin GPU dedicada). Este apartado documenta las tres fases de ese proceso y la evidencia experimental que condujo a la arquitectura mixta finalmente desplegada.

\textbf{Línea base: Qwen2.5 (abril--mayo 2026).} El proyecto arrancó con una configuración de dos modelos de la familia Qwen2.5~\cite{hui_qwen25_coder_2024}: \texttt{qwen2.5-coder:1.5b} como clasificador ligero (Tier~1, 300--500~ms, JSON determinista) y \texttt{qwen2.5-coder:7b} como razonador (Tier~2, 5--30~s). Esta configuración estableció un estándar de rendimiento estable y predecible que cumplía holgadamente el requisito RNF-03 (latencia inferior a 2~s para comandos directos).\\[2pt]
El modelo \texttt{bge-m3}~\cite{bge_m3} se mantuvo desde el inicio para \textit{embeddings} semánticos por su eficiencia (500~MB en RAM) y capacidad multilingüe.

\textbf{Iteración experimental: Qwen3.5 (mayo--junio 2026).} En mayo de 2026 se evaluó la familia Qwen3.5~\cite{qwen3_5} como posible sustituta. El atractivo era claro: ventana de contexto ampliada a 256K tokens, soporte para 201 idiomas y modo \textit{thinking} nativo.\\[2pt]
Se desplegaron los modelos \texttt{qwen3.5:2b}, \texttt{qwen3.5:4b} y \texttt{qwen3.5:9b} en Ollama y se sometieron a pruebas de estrés con el corpus de comandos reales del sistema. Los resultados fueron concluyentes:

\begin{itemize}
    \item \textbf{Latencia en CPU:} el modo \textit{thinking} de Qwen3.5 es forzado y no desactivable vía API. En el Ryzen 5, la generación de tokens de razonamiento previos a la respuesta producía una tasa de 0.6--1.2~tok/s. El modelo de 9B alcanzaba latencias de 3 a 12 minutos por consulta conversacional, invalidando cualquier aplicación en tiempo real.
    \item \textbf{Inestabilidad en JSON:} a diferencia de su predecesor, Qwen3.5 mostró inconsistencias en la generación de JSON estructurado con \texttt{temperature=0.0}, produciendo ocasionalmente campos ausentes o claves mal formadas que rompían el \textit{parser} del enrutador.
    \item \textbf{Consumo de RAM:} el modelo de 9B elevaba la RAM total del sistema a 7.9~GB, un 18\% más que la configuración base, dejando menos margen para picos de carga.
\end{itemize}

La conclusión experimental fue que Qwen3.5 está diseñado para entornos con GPU y su arquitectura de \textit{thinking} forzado lo hace inviable en CPU para aplicaciones con requisitos de latencia.

\textbf{Arquitectura mixta final (junio 2026).} A la vista de los resultados, se adoptó una estrategia diferenciada por niveles: mantener el modelo óptimo para cada tarea en lugar de usar una familia única para todo el sistema.

\begin{itemize}
    \item \textbf{Tier~1 --- Clasificación:} se conservó \texttt{qwen2.5-coder:1.5b}. Su generación de JSON determinista, su latencia de 300--500~ms y su consumo de solo 986~MB en RAM lo hacen insustituible para el 92\% de los comandos que requieren clasificación. Es un modelo maduro, predecible y perfectamente adaptado al hardware disponible.
    \item \textbf{Tier~2 --- Razonamiento:} se seleccionó \texttt{mistral:7b-instruct-q4\_K\_M}~\cite{mistral_7b}. Frente a Qwen3.5, Mistral ofrece respuesta directa sin \textit{thinking} forzado, con latencias de 27--60~s en CPU y una calidad conversacional que evita la sensación de estar hablando con un \textit{bot} de reglas. La cuantización Q4\_K\_M reduce su huella a 4.4~GB, permitiendo la coexistencia con el clasificador de Tier~1 y \texttt{bge-m3} dentro de los 16~GB del servidor.
\end{itemize}

La Tabla~\ref{tab:modelos_evaluados} resume la comparativa de latencias observadas durante las pruebas de estrés. La Figura~\ref{fig:latencia_comparativa} del Capítulo~9 muestra el impacto visual de estas diferencias.

\begin{table}[ht]
\centering
\small
\setlength{\tabcolsep}{5pt}
\begin{tabular}{|p{2.5cm}|p{3.2cm}|p{2.0cm}|p{2.0cm}|p{2.0cm}|}
\hline
\textbf{Modelo} & \textbf{Tarea} & \textbf{Latencia (s)} & \textbf{Fallos JSON} & \textbf{RAM (GB)} \\
\hline
Qwen2.5 1.5B & Clasificación (Tier~1) & 0.3--0.5 & 0\% & 0.99 \\
Qwen2.5 7B & Razonamiento (Tier~2) & 5--30 & 0\% & 6.0 \\
Qwen3.5 4B & Razonamiento (eval.) & 45--180 & 12\% & 3.8 \\
Qwen3.5 9B & Razonamiento (eval.) & 180--720 & 18\% & 7.1 \\
Mistral 7B (Q4) & Razonamiento (Tier~2) & 27--60 & 0\% & 4.4 \\
\hline
\end{tabular}
\caption{Comparativa de latencias y fiabilidad de los modelos evaluados para el Tier~2. Las pruebas se realizaron sobre el corpus de 500 comandos reales del proyecto en un Ryzen 5 con 16~GB RAM.}
\label{tab:modelos_evaluados}
\end{table}

Esta evolución ilustra una lección central del proyecto: un TFG de ingeniería no es un camino lineal de adopción tecnológica, sino un proceso de diagnóstico basado en evidencia. La decisión de descartar Qwen3.5 tras evaluarlo y optar por una arquitectura mixta ---cada modelo especializado en su nivel--- demuestra un enfoque guiado por el cumplimiento estricto de los requisitos no funcionales frente a la mera adopción de tecnologías de tendencia. El modelo \texttt{bge-m3} se mantuvo desde el inicio por su eficiencia y su capacidad multilingüe, pues el autor trabaja tanto en español como en inglés.

\subsection{Aislamiento de Red: La Arquitectura Air-Gapped}

Durante las primeras pruebas se observó que los dispositivos comerciales (Xiaomi, Tuya) intentaban abrir canales salientes hacia sus respectivas nubes en cuanto se conectaban a la red doméstica. Era necesario impedir esa comunicación.

La solución fue desplegar una arquitectura de red \textbf{100\% air-gapped}: un segmento \texttt{192.168.1.0/24} sobre un router TP-Link TL-MR6400 sin tarjeta SIM (la ranura está físicamente rota) y sin ningún tipo de enlace WAN. Este router actúa exclusivamente como \textit{switch} Ethernet y punto de acceso WLAN. Todos los dispositivos del sistema (servidor, portátil de desarrollo, enchufe Tuya, sensores BLE) conviven en este segmento aislado. El aislamiento físico es la única garantía verificable de que ningún paquete generado por los dispositivos comerciales pueda alcanzar Internet.

\subsubsection{Hotspot Spoofing para Dispositivos Tuya}

El aislamiento físico introduce una fricción operativa: los dispositivos Tuya necesitan conectarse a Internet durante el emparejamiento inicial para validar el registro contra la nube del fabricante. Para resolverlo sin renunciar al air-gapped, se aplica la técnica de \textit{Hotspot Spoofing}:

\begin{enumerate}
    \item Se levanta temporalmente un punto de acceso móvil (teléfono Android) que replica el SSID y la contraseña del router aislado.
    \item Se completa el emparejamiento del dispositivo Tuya contra la nube del fabricante.
    \item Durante el emparejamiento, se extrae la \textit{Local Key} del dispositivo mediante \texttt{tinytuya wizard}.
    \item Se apaga el \textit{hotspot}. El dispositivo, al encontrar el SSID original (el router TP-Link), se conecta al segmento air-gapped.
    \item A partir de ese momento, todas las órdenes viajan directamente a la IP local del dispositivo (\texttt{192.168.1.100}) mediante el protocolo 3.5 de \texttt{tinytuya}, sin volver a depender de la infraestructura de Tuya.
\end{enumerate}

La \textit{Local Key} queda almacenada exclusivamente en el servidor local y nunca se transmite a terceros.

\subsubsection{Lectura de Sensores Xiaomi LYWSD03MMC}

Los sensores de temperatura Xiaomi presentaban un problema similar: el firmware de fábrica cifra los anuncios BLE con una \textit{bind key} vinculada a la cuenta del fabricante, lo que reintroduce dependencia de la nube. La ruta prevista era flashear el firmware alternativo ATC MiThermometer~\cite{pvvx_atc_mithermometer}, pero durante las pruebas de campo se descubrió una alternativa mejor.

Al conectar al sensor mediante \texttt{bleak}~\cite{bleak} y obtener un volcado de sus servicios GATT, se descubrió que la característica \texttt{ebe0ccc1} del servicio propietario de Xiaomi es legible \textbf{sin autenticación}. Devuelve 5 bytes con temperatura ($\times$~100), humedad relativa y voltaje de batería:

\begin{itemize}
    \item Bytes 0--1: temperatura en \texttt{int16} LE, con signo. Ejemplo: \texttt{0x09cb} = 2507 $\rightarrow$ 25.07$^\circ$C.
    \item Byte 2: humedad relativa (\%). Ejemplo: \texttt{0x3b} = 59\%.
    \item Bytes 3--4: voltaje de batería en mV (\texttt{uint16} LE). Ejemplo: \texttt{0x0bef} = 3055 mV.
\end{itemize}

Esta vía se adoptó como método principal de integración. El flasheo ATC se mantiene como alternativa para escenarios futuros que requieran lectura pasiva de múltiples sensores simultáneos. La decisión preserva la integridad del hardware (sin riesgo de \textit{brick}) y elimina cualquier dependencia de la nube de Xiaomi. El \textit{trade-off} es que la conexión GATT consume unos 12 segundos por ciclo de lectura, durante los cuales el sensor no puede ser leído por otro dispositivo. En un escenario con un único punto de recolección, esta limitación es aceptable.

\subsection{Selección del Sistema Operativo y el Servidor}

La inferencia de modelos de lenguaje mediante Ollama demanda recursos que no pueden coexistir de forma estable con un entorno de desarrollo interactivo en el mismo hardware. Por eso se adoptó una arquitectura de dos nodos físicos:

\begin{itemize}
    \item \textbf{Servidor dedicado:} torre con APU Ryzen y 16 GB de RAM, ejecutando \textbf{Ubuntu Server 24.04 LTS} en modo \textit{headless} (sin interfaz gráfica). Alberga el \textit{runtime} completo del sistema: Ollama, OpenClaw, los \textit{drivers} Python de dispositivos y la base de datos SQLite. La ausencia de escritorio libera entre 500 y 800 MB de RAM que se destinan a los modelos de lenguaje.
    \item \textbf{Portátil de desarrollo:} conserva las herramientas de ingeniería (compilación de firmware, redacción \LaTeX, Git) y se conecta al servidor exclusivamente por SSH con llaves RSA. Esta separación garantiza que los picos de inferencia no compitan con las tareas interactivas del desarrollador, y viceversa.
\end{itemize}

\subsubsection{Por qué Ubuntu Server 24.04 LTS}

Se eligió Ubuntu Server por su soporte nativo para todas las dependencias del proyecto (Python 3.12+, Node.js 24, \texttt{bluez} para BLE, \texttt{systemd}) y su ciclo de soporte extendido de 5 años. Dentro de Ubuntu, se optó por la versión 24.04 LTS frente a la recién publicada 26.04 LTS por tres razones:

\begin{enumerate}
    \item \textbf{Madurez del ecosistema:} Ollama, Python, \texttt{bluez} y \texttt{systemd} llevan más de dos años verificados en 24.04. La versión 26.04 introduce versiones renovadas cuyo comportamiento en entornos air-gapped no está suficientemente probado.
    \item \textbf{Disponibilidad de \textit{wheels} binarios:} las librerías Python con dependencias nativas (\texttt{bleak}, \texttt{tinytuya}) tienen paquetes precompilados estables para 24.04. En 26.04 podrían requerir compilar desde el código fuente, algo inviable sin acceso a Internet.
    \item \textbf{Estabilidad:} los primeros meses de cualquier LTS concentran la mayor densidad de errores críticos. Operar sobre 24.04 reduce drásticamente el riesgo de encontrar un fallo en un componente del sistema que, en un entorno air-gapped, requeriría intervención física.
\end{enumerate}

Se consideraron otras opciones que fueron descartadas. \textbf{Linux Mint}, incluso en su variante ligera XFCE, consume entre 400 y 700 MB de RAM en reposo debido a los servicios gráficos. Es memoria que en un servidor \textit{headless} se destina íntegramente a los modelos de lenguaje. Las distribuciones \textit{rolling release} como \textbf{CachyOS} o \textbf{Arch Linux} presentan un riesgo adicional en un entorno aislado: una actualización del kernel o de \texttt{bluez} podría romper la compatibilidad con el adaptador Bluetooth, y al no tener acceso a Internet para resolverlo, el sistema quedaría inoperativo hasta una intervención física. Ambos enfoques resultaron incompatibles con los requisitos de estabilidad del proyecto.

\subsubsection{Herramientas de Gestión}

Para la administración cotidiana del servidor se despliega \textbf{Cockpit} (puerto 9090), una interfaz web ligera que permite monitorizar servicios, CPU y RAM sin necesidad de acceso físico. El desarrollo se realiza mediante \textbf{SSH con autenticación por par de llaves RSA} desde el portátil, lo que permite despliegue remoto de código (\texttt{git pull}) y administración de servicios sin interfaz gráfica.


\section{Metodología de Trabajo}

\subsection{Control de Versiones y Propiedad Intelectual}

El código fuente de Luxe Core AI no es solo un entregable académico: representa la base sobre la que podría construirse una futura iniciativa empresarial. Por eso se aloja en un repositorio privado de GitHub con control de versiones Git. Esta infraestructura no actúa únicamente como mecanismo de \textit{backup} o \textit{rollback}, sino que audita cada contribución y restringe el acceso público al núcleo del proyecto. La trazabilidad que ofrece Git permite, además, vincular cada decisión de diseño documentada en esta memoria con el código que la implementa.

\subsection{Despliegue Inicial: Incidencias y Soluciones}

La puesta en marcha del servidor en la red air-gapped enfrentó dos obstáculos que ilustran la clase de problemas que surgen al trabajar con hardware real.

\subsubsection{El Puerto LAN/WAN Híbrido}

Al conectar el servidor por Ethernet al router TP-Link, el portátil (conectado por Wi-Fi al mismo router) devolvía un error de enrutamiento al intentar acceder por SSH. El cable se había conectado por error al puerto híbrido LAN/WAN del router. Este puerto está diseñado para conectarse a un módem externo, y la lógica interna del \textit{firmware} del TP-Link, al no detectar una conexión WAN válida, mantiene la interfaz en un estado de aislamiento eléctrico que impide el intercambio de tramas Ethernet a nivel de Capa~2 con el resto del \textit{switch} interno.

La solución fue migrar el cable a un puerto LAN dedicado. El \textit{switch} interno restableció inmediatamente el intercambio de tramas, y el servidor apareció en la red con su IP estática \texttt{192.168.1.50}.

\subsubsection{Pasarela de Red de Contingencia}

La instalación mínima de Ubuntu Server 24.04 LTS carece de herramientas de edición de texto y del demonio \texttt{wpasupplicant} para gestionar la antena Wi-Fi. Sin conectividad a Internet en el segmento air-gapped, el servidor no podía instalar las dependencias del proyecto.

Se implementó una solución temporal utilizando el portátil de desarrollo como \textbf{pasarela NAT} entre el segmento air-gapped y una conexión celular externa (anclaje USB desde un smartphone). La configuración se realizó en dos frentes.

\textbf{En el servidor (Netplan):} se estableció la IP estática y se definió al portátil como \textit{default gateway}:

\begin{verbatim}
network:
  version: 2
  renderer: networkd
  ethernets:
    enp3s0:
      addresses: [192.168.1.50/24]
      routes:
        - to: default
          via: 192.168.1.101
      nameservers:
        addresses: [8.8.8.8, 8.8.4.4]
\end{verbatim}

\textbf{En el portátil:} se ejecutaron los comandos para habilitar el reenvío de paquetes y el enmascaramiento NAT:

\begin{verbatim}
sudo sysctl -w net.ipv4.ip_forward=1
sudo iptables -t nat -A POSTROUTING \
  -o enp0s20f0u2u2 -j MASQUERADE
sudo iptables -P FORWARD ACCEPT
\end{verbatim}

La reconfiguración de la cadena \texttt{FORWARD} fue necesaria porque la política por defecto de \texttt{iptables} en las estaciones de trabajo es \texttt{DROP}: cualquier paquete que intente cursarse entre interfaces distintas del mismo \textit{host} es descartado. Durante la ventana de contingencia, reducir esta protección es seguro porque el segmento air-gapped no tiene salida a Internet en operación normal.

Con esta pasarela, el servidor pudo ejecutar \texttt{apt update \&\& apt dist-upgrade} e instalar todo el \textit{stack} de desarrollo: Python 3.12+, Node.js 24, Git, Cockpit y \texttt{bluez}/\texttt{btmgmt}. Una vez completada la instalación, la política \texttt{FORWARD} se restauró a \texttt{DROP} y el \textit{forwarding} se desactivó (\texttt{net.ipv4.ip\_forward=0}), devolviendo el segmento a su estado de aislamiento total.

\subsection{Rigor Documental}

Toda la documentación del proyecto (esta memoria, las decisiones de diseño, los diagramas) se redacta en \LaTeX{} y se aloja en el mismo repositorio Git que el código fuente. Este enfoque de \textit{Docs as Code} asegura trazabilidad entre las decisiones de diseño, el código que las implementa y la documentación que las describe, además de proporcionar consistencia tipográfica profesional y control de versiones unificado.

\subsection{Uso de la Inteligencia Artificial en el Desarrollo}

Durante el desarrollo de este proyecto se ha utilizado inteligencia artificial como herramienta de apoyo, tanto para la generación de fragmentos de código como para la redacción de partes de esta memoria, siempre bajo criterio de supervisión humana directa. Las herramientas se eligieron buscando la mejor relación calidad-precio disponible en el momento del TFG ---sin pretender predecir el futuro del sector--- e incluyen:

\begin{itemize}
    \item \textbf{OpenCode}~\cite{opencode}: asistente de codificación interactivo, software libre, ejecutado directamente en el servidor de desarrollo a través de SSH. Se utilizó para depuración, refactorización y generación de fragmentos de código Python y Bash, así como para la edición y corrección de esta memoria en \LaTeX{}.

    \item \textbf{DeepSeek V4}~\cite{deepseek_api}: utilizado como modelo de inferencia desde OpenCode a través de su API, no como chat web. En el momento del desarrollo ofrecía la mejor relación entre capacidad de razonamiento y coste por token de los proveedores cloud disponibles. No forma parte del sistema en producción, que utiliza exclusivamente modelos locales a través de Ollama.

    \item \textbf{Otros modelos} (ChatGPT, Gemini): utilizados puntualmente para consultas específicas durante la fase de investigación inicial, cuando el proyecto aún no disponía de infraestructura local de inferencia.
\end{itemize}

Cada línea de código y cada párrafo generado han sido revisados, comprendidos, validados y, cuando ha sido necesario, modificados por el autor. No hay nada en este trabajo que no haya pasado por mi criterio personal. En esta forma de trabajar, la inteligencia artificial ha actuado como asistente de propuestas, nunca como autoridad incuestionable.

Definidas las restricciones y la forma de trabajo, el capítulo siguiente detalla los requisitos funcionales y no funcionales que debe cumplir el sistema.
